% chktex-file 19
% chktex-file 35
% chktex-file 36
 \documentclass[12pt]{article}
 \usepackage[left=2cm,right=2cm,top=2cm,bottom=2cm,bindingoffset=0cm]{geometry}
 \usepackage{amssymb}
 \usepackage{amsmath}
 \usepackage{amsthm}
 \usepackage{enumerate}
 \usepackage[T1,T2A]{fontenc}
 \usepackage[utf8]{inputenc}
 \usepackage[matrix,arrow,curve]{xy}
 \usepackage[colorlinks=true, urlcolor=blue, linkcolor=blue, citecolor=blue,
     pdfborder={0 0 0}]{hyperref}
\usepackage{enumitem}

 %------------------------------------------------------------------------------
 \newtheorem{theorem}{Theorem}[section]
 \newtheorem{lemma}[theorem]{Lemma}
 \newtheorem{proposition}[theorem]{Proposition}
 \newtheorem{remark}[theorem]{Remark}
 \newtheorem{corollary}[theorem]{Corollary}
 \newtheorem{definition}[theorem]{Definition}
 \newtheorem{example}[theorem]{Example}
 \newtheorem{question}[theorem]{Question}

 \newcommand{\projtens}{\mathbin{\widehat{\otimes}}}
 \newcommand{\isom}[1]{\mathop{\mathbin{\cong}}\limits_{#1}}
 \newcommand{\Nor}{\mathbf{Nor}}
 \newcommand{\Ban}{\mathbf{Ban}}
 \newcommand{\lo}[1]{\ell_1^0(#1)}
 \newcommand{\ellone}[1]{\ell_1(#1)}
 \newcommand{\norm}[1]{\lVert #1\rVert}
 \newcommand{\supp}{\operatorname{supp}}
 \newcommand{\ext}{\operatorname{ext}}
 \newcommand{\conv}{\operatorname{conv}}
 \newcommand{\Lin}{\operatorname{span}}
 \newcommand{\Im}{\operatorname{Im}}
 \newcommand{\Ker}{\operatorname{Ker}}
 \newcommand{\BM}{d_{BM}}
 %------------------------------------------------------------------------------

\begin{document}

\begin{center}
    \Large \textbf{Extremally projective normed spaces}\\[0.5cm]
    \small {N. T. Nemesh}\\[0.5cm]
\end{center}

\thispagestyle{empty}

\medskip
\textbf{Abstract:} A normed space is called extremally projective if it is
$C$-topologically projective in the category of normed spaces for every $C>1$.
Helemskii proved that the spaces $\lo{\Lambda}$ of finitely supported functions
on a set $\Lambda$ with the $\ell_1$-norm are extremally projective and asked
whether every extremally projective normed space is isometrically isomorphic to
some $\lo{\Lambda}$. We answer this question in the negative and show that
extremal projectivity is an almost isometric rather than an isometric property:
a normed space $P$ is extremally projective if and only if for every
$\varepsilon>0$ the Banach-Mazur distance between $P$ and $\lo{\Lambda}$ is less
than $1+\varepsilon$, where $\Lambda$ is any set of cardinality equal to the
algebraic dimension of $P$. Along the way we show that extremally projective
normed spaces contain no infinite dimensional complete subspaces, so that no
infinite dimensional Banach space is extremally projective in the category of
normed spaces.
\medskip

\textbf{Keywords:} projective normed space, topological projectivity, extremal
projectivity, Banach-Mazur distance, $\ell_1$-space.

\bigskip

%-------------------------------------------------------------------------------
%	Introduction
%-------------------------------------------------------------------------------

\section{Introduction}\label{SectionIntroduction}

Projectivity in categories of normed and Banach modules admits several
non-equivalent versions. Besides the classical relative projectivity of
Helemskii~\cite{HelHomolBanTopAlg}, which ignores the norms of lifting
operators, there are so-called metric and topological
projectivity~\cite{HelMetrFrQMod, WhiteInjmoduAlg, NemMetrTopHomol}, where the
norm of a lifting is controlled by a constant. In the metric theory this
constant is $1$, and in the topological theory it is some fixed $C\geq 1$. The
present paper is devoted to the limiting case between these two theories: a
module is called extremally projective if it is $C$-topologically projective for
every $C>1$. This notion was introduced and studied by Helemskii under the name
of extreme projectivity~\cite{HelExtrProjSeqAlg, HelMetrFrQMod}. We restrict
ourselves to the simplest possible algebra, namely the zero algebra. Normed
modules over the zero algebra are just normed spaces, so we deal with extremal
projectivity in the category $\Nor$ of normed spaces and bounded linear
operators.

For Banach spaces the corresponding questions are classical. By a theorem of
K\"othe~\cite{KotheTopProjBanSp} topologically projective Banach spaces are
exactly the Banach spaces topologically isomorphic to $\ellone{\Lambda}$ for some
set $\Lambda$, and by a theorem of Grothendieck~\cite{GrothMetrProjFlatBanSp}
a Banach space has the metric lifting property if and only if it is
isometrically isomorphic to some $\ellone{\Lambda}$. For normed spaces the
answers are different, because the category $\Nor$ has its own free objects.
These are the spaces $\lo{\Lambda}$ of finitely supported functions on a set
$\Lambda$ with the $\ell_1$-norm. Gr{\o}nb{\ae}k~\cite{GronLiftProbNormSp}
extended K\"othe's theorem to $\Nor$: a normed space is topologically projective
if and only if it is topologically isomorphic to $\lo{\Lambda}$ for some set
$\Lambda$. Helemskii~\cite{HelMetrFrQMod} proved the metric counterpart: a
normed space is metrically projective if and only if it is isometrically
isomorphic to some $\lo{\Lambda}$. He also showed that the spaces $\lo{\Lambda}$
are extremally projective and asked whether, conversely, every extremally
projective normed space is isometrically isomorphic to some $\lo{\Lambda}$
\cite[section 2]{HelExtrProjSeqAlg}, \cite[section 1]{HelMetrFrQMod}.

The main results of this paper are the following. First, the answer to
Helemskii's question is negative: the subspace of $\ell_1$ spanned by the
finitely supported sequences and a single sequence with infinite support is
extremally projective but not isometrically isomorphic to any $\lo{\Lambda}$.
Second, extremal projectivity turns out to be an almost isometric property: a
normed space $P$ is extremally projective if and only if for every
$\varepsilon>0$ there is a linear bijection $T:P\to\lo{\Lambda}$ with
$\norm{T}\norm{T^{-1}}<1+\varepsilon$, where $\Lambda$ is a set of cardinality
$\dim P$. Third, extremally projective normed spaces have no infinite
dimensional complete subspaces. In particular, an infinite dimensional Banach
space is never extremally projective in $\Nor$, although it may well be
extremally projective in the category $\Ban$ of Banach spaces. The last fact
explains the main technical difficulty of the subject: when the lifting
problems are posed for all normed spaces rather than Banach spaces, the
standard tools of Banach space geometry, which rely on completeness, are no
longer available, and we must work with the finite dimensional pieces of our
spaces directly.

Let us fix the notation and recall the definitions. All normed spaces are
considered over the complex field; the real case is identical. For a normed
space $E$ we denote by $B_E$ its closed unit ball and by $B_E^\circ$ its open
unit ball. By $\Nor$ we denote the category of normed spaces and bounded linear
operators and by $\Ban$ its full subcategory of Banach spaces. For a set
$\Lambda$ we denote by $\ellone{\Lambda}$ the Banach space of absolutely
summable families and by $\lo{\Lambda}$ its dense subspace consisting of finitely
supported families. For $\lambda\in\Lambda$ the symbol $e_\lambda$ denotes the
family which equals $1$ at $\lambda$ and $0$ elsewhere; the family
${(e_\lambda)}_{\lambda\in\Lambda}$ is a Hamel basis of $\lo{\Lambda}$. For a
normed space $E$ and a set $\Lambda$ there is a bijection between bounded
families ${(x_\lambda)}_{\lambda\in\Lambda}$ in $E$ and bounded linear
operators $T:\lo{\Lambda}\to E$, given by $T(e_\lambda)=x_\lambda$, and
$\norm{T}=\sup_{\lambda\in\Lambda}\norm{x_\lambda}$. This is the sense in
which the spaces $\lo{\Lambda}$ are free objects of $\Nor$; note that no
completeness of $E$ is required here, in contrast with the case of
$\ellone{\Lambda}$.

Following~\cite{NemMetrTopHomol}, we say that a bounded linear operator
$\xi:X\to Y$ between normed spaces is $c$-topologically surjective, where
$c\geq 1$, if $c\,\xi(B_X^\circ)\supset B_Y^\circ$, and that $\xi$ is
coisometric if it is contractive and $1$-topologically surjective, that is, if
$\xi(B_X^\circ)=B_Y^\circ$. Equivalently, $\xi$ is $c$-topologically surjective
if for every $y\in Y$ and every $\delta>0$ there exists $x\in X$ with
$\xi(x)=y$ and $\norm{x}<c\norm{y}+\delta$. A linear bijection $T:E\to F$ is
called a $c$-isomorphism if $\norm{T}\norm{T^{-1}}\leq c$. The Banach-Mazur
distance between $E$ and $F$ is
\[
\BM(E,F)=\inf\{\norm{T}\norm{T^{-1}}:T:E\to F \mbox{ is a linear bijection}\},
\]
with $\BM(E,F)=+\infty$ if there are no bounded linear bijections with bounded
inverses between $E$ and $F$. A bounded linear operator $\zeta:E\to F$ is
called a $c$-retraction if there exists a bounded linear operator $\eta:F\to E$
with $\zeta\eta=1_F$ and $\norm{\zeta}\norm{\eta}\leq c$. In this case $F$ is
called a $c$-retract of $E$.

\begin{definition}\label{DefTopProj} Let $C\geq 1$. A normed space $P$ is
    called $C$-topologically projective if for every $c\geq 1$, every
    $c$-topologically surjective operator $\xi:X\to Y$ between normed spaces
    and every bounded linear operator $\phi:P\to Y$ there exists a bounded
    linear operator $\psi:P\to X$ such that $\xi\psi=\phi$ and
    $\norm{\psi}\leq cC\norm{\phi}$. A normed space is called topologically
    projective if it is $C$-topologically projective for some $C\geq 1$, and
    extremally projective if it is $C$-topologically projective for every
    $C>1$.
\end{definition}

The operator $\psi$ in definition~\ref{DefTopProj} is called a lifting of
$\phi$ along $\xi$. Extremal projectivity is precisely the extreme projectivity
of Helemskii, see proposition~\ref{EquivWithHelemskii} below. We emphasize that
the spaces $X$ and $Y$ in definition~\ref{DefTopProj} are arbitrary normed
spaces. If we restrict them to Banach spaces, we obtain the notions of
topological and extremal projectivity in $\Ban$, which we shall mention only
for comparison. The two notions differ substantially, see
section~\ref{SectionNecessaryConditions}.

The paper is organized as follows. In section~\ref{SectionBasicTheory} we
collect the basic properties of extremally projective normed spaces: the
spaces $\lo{\Lambda}$ are extremally projective, extremal projectivity passes
to $(1+\varepsilon)$-retracts, and a normed space is extremally projective if
and only if it is a $(1+\varepsilon)$-retract of a free normed space for every
$\varepsilon>0$. In section~\ref{SectionNecessaryConditions} we show that
topologically projective normed spaces do not contain infinite dimensional
Banach subspaces. Section~\ref{SectionCounterexample} contains the
counterexample to Helemskii's question. Section~\ref{SectionMainTheorem} is
devoted to the almost isometric characterization of extremally projective
normed spaces.

\end{document}
